\documentclass[../DoAn.tex]{subfiles}
\begin{document}

\section*{Phụ lục: Một số ảnh spectrogram minh họa}
\addcontentsline{toc}{section}{Phụ lục: Một số ảnh spectrogram minh họa}
\label{sec:appendix_spectrogram_examples}

Phụ lục này bổ sung một số ảnh spectrogram để minh họa trực quan cho dữ liệu và các kịch bản đánh giá trong đồ án.
Các ảnh được đặt tại \texttt{Hinh\_ve/phu\_luc/anh\_phu\_luc}; mỗi nhóm hình tương ứng với một loại dữ liệu và các
ảnh trong cùng nhóm được xếp trên một hàng ngang.

\newcommand{\appendixpic}[3]{%
    \begin{minipage}{#1}
        \centering
        \includegraphics[width=\linewidth]{#2}
        {\scriptsize #3}
    \end{minipage}%
}

\subsection*{Mẫu từ tập DroneDetectv2}

Hình~\ref{fig:appendix_dataset_drone} minh họa hai mẫu spectrogram thuộc nguồn dữ liệu drone chính được sử dụng trong
tập huấn luyện.

\begin{figure}[H]
    \centering
    \appendixpic{0.30\textwidth}{Hinh_ve/phu_luc/anh_phu_luc/dataset_drone/1.png}{Mẫu 1}
    \hspace{0.08\textwidth}
    \appendixpic{0.30\textwidth}{Hinh_ve/phu_luc/anh_phu_luc/dataset_drone/2.png}{Mẫu 2}
    \caption{Một số mẫu spectrogram từ tập DroneDetectv2}
    \label{fig:appendix_dataset_drone}
\end{figure}

\subsection*{Mẫu từ các bộ dữ liệu ngoài miền}

Hình~\ref{fig:appendix_dataset_external} trình bày các mẫu từ RFUAV và DroneRF, được dùng để kiểm tra khả năng tổng
quát hóa của mô hình trên nguồn dữ liệu khác với tập huấn luyện.

\begin{figure}[H]
    \centering
    \appendixpic{0.30\textwidth}{Hinh_ve/phu_luc/anh_phu_luc/dataset_ngoai/RFUAV.png}{RFUAV}
    \hspace{0.08\textwidth}
    \appendixpic{0.30\textwidth}{Hinh_ve/phu_luc/anh_phu_luc/dataset_ngoai/DroneRF.png}{DroneRF}
    \caption{Một số mẫu spectrogram từ các bộ dữ liệu ngoài miền}
    \label{fig:appendix_dataset_external}
\end{figure}

\subsection*{Mẫu dữ liệu RF tự thu}

Hình~\ref{fig:appendix_self_collected} trình bày bốn mẫu spectrogram sinh ra từ dữ liệu RF tự thu. Nhóm ảnh này phản
ánh điều kiện thu bằng SDR trong môi trường thực tế hơn so với dữ liệu công khai dùng để huấn luyện.

\begin{figure}[H]
    \centering
    \appendixpic{0.23\textwidth}{Hinh_ve/phu_luc/anh_phu_luc/tu_thu/1.png}{Mẫu 1}
    \hfill
    \appendixpic{0.23\textwidth}{Hinh_ve/phu_luc/anh_phu_luc/tu_thu/2.png}{Mẫu 2}
    \hfill
    \appendixpic{0.23\textwidth}{Hinh_ve/phu_luc/anh_phu_luc/tu_thu/3.png}{Mẫu 3}
    \hfill
    \appendixpic{0.23\textwidth}{Hinh_ve/phu_luc/anh_phu_luc/tu_thu/4.png}{Mẫu 4}
    \caption{Một số mẫu spectrogram từ dữ liệu RF tự thu}
    \label{fig:appendix_self_collected}
\end{figure}

\subsection*{Mẫu thu trong nhà và ngoài trời}

Hình~\ref{fig:appendix_indoor} và Hình~\ref{fig:appendix_outdoor} minh họa hai điều kiện thu khác nhau. Môi trường thu,
vật cản, phản xạ đa đường và mức nhiễu nền có thể làm thay đổi hình dạng năng lượng trên spectrogram.

\begin{figure}[H]
    \centering
    \appendixpic{0.32\textwidth}{Hinh_ve/phu_luc/anh_phu_luc/indoor/1.png}{Indoor}
    \caption{Mẫu spectrogram thu trong nhà}
    \label{fig:appendix_indoor}
\end{figure}

\begin{figure}[H]
    \centering
    \appendixpic{0.32\textwidth}{Hinh_ve/phu_luc/anh_phu_luc/outdoor/1.png}{Outdoor}
    \caption{Mẫu spectrogram thu ngoài trời}
    \label{fig:appendix_outdoor}
\end{figure}

\subsection*{Mẫu chất lượng tốt và chất lượng kém}

Hình~\ref{fig:appendix_good_examples} minh họa các mẫu spectrogram có cấu trúc tín hiệu tương đối rõ. Ngược lại,
Hình~\ref{fig:appendix_not_good_examples} trình bày các mẫu khó quan sát hơn, có thể do SNR thấp, nhiễu đồng kênh,
clipping hoặc điều kiện thu không ổn định.

\begin{figure}[H]
    \centering
    \appendixpic{0.30\textwidth}{Hinh_ve/phu_luc/anh_phu_luc/good/1.png}{Mẫu 1}
    \hfill
    \appendixpic{0.30\textwidth}{Hinh_ve/phu_luc/anh_phu_luc/good/2.png}{Mẫu 2}
    \hfill
    \appendixpic{0.30\textwidth}{Hinh_ve/phu_luc/anh_phu_luc/good/3.png}{Mẫu 3}
    \caption{Một số mẫu spectrogram có cấu trúc tín hiệu rõ}
    \label{fig:appendix_good_examples}
\end{figure}

\begin{figure}[H]
    \centering
    \appendixpic{0.30\textwidth}{Hinh_ve/phu_luc/anh_phu_luc/not_good/1.png}{Mẫu 1}
    \hfill
    \appendixpic{0.30\textwidth}{Hinh_ve/phu_luc/anh_phu_luc/not_good/2.png}{Mẫu 2}
    \hfill
    \appendixpic{0.30\textwidth}{Hinh_ve/phu_luc/anh_phu_luc/not_good/3.png}{Mẫu 3}
    \caption{Một số mẫu spectrogram khó quan sát hơn}
    \label{fig:appendix_not_good_examples}
\end{figure}

\subsection*{Mẫu tín hiệu nền và nhiễu cùng băng tần}

Hình~\ref{fig:appendix_noise_examples} minh họa một số mẫu tín hiệu nền hoặc nguồn phát cùng băng tần. Các nguồn này
có thể chồng lấn với tín hiệu drone trong băng 2,4~GHz và làm tăng độ khó của bài toán phân loại.

\begin{figure}[H]
    \centering
    \appendixpic{0.30\textwidth}{Hinh_ve/phu_luc/anh_phu_luc/nhieu/blue.png}{Bluetooth}
    \hfill
    \appendixpic{0.30\textwidth}{Hinh_ve/phu_luc/anh_phu_luc/nhieu/wf1.png}{Wi-Fi -- 1}
    \hfill
    \appendixpic{0.30\textwidth}{Hinh_ve/phu_luc/anh_phu_luc/nhieu/wf2.png}{Wi-Fi -- 2}
    \caption{Một số mẫu tín hiệu nền và nguồn phát cùng băng tần}
    \label{fig:appendix_noise_examples}
\end{figure}

\subsection*{Mẫu spectrogram theo khoảng cách thu}

Hình~\ref{fig:appendix_distance_examples} minh họa các mẫu spectrogram ở những khoảng cách thu khác nhau. Khoảng cách
có thể ảnh hưởng đến công suất thu và SNR, nhưng mức thay đổi trên spectrogram còn phụ thuộc vào hướng anten, môi
trường truyền sóng và nhiễu nền tại thời điểm thu.

\begin{figure}[H]
    \centering
    \appendixpic{0.155\textwidth}{Hinh_ve/phu_luc/anh_phu_luc/khoang_cach/10m.png}{10~m}
    \hfill
    \appendixpic{0.155\textwidth}{Hinh_ve/phu_luc/anh_phu_luc/khoang_cach/30m.png}{30~m}
    \hfill
    \appendixpic{0.155\textwidth}{Hinh_ve/phu_luc/anh_phu_luc/khoang_cach/40m.png}{40~m}
    \hfill
    \appendixpic{0.155\textwidth}{Hinh_ve/phu_luc/anh_phu_luc/khoang_cach/60m.png}{60~m}
    \hfill
    \appendixpic{0.155\textwidth}{Hinh_ve/phu_luc/anh_phu_luc/khoang_cach/80m.png}{80~m}
    \hfill
    \appendixpic{0.155\textwidth}{Hinh_ve/phu_luc/anh_phu_luc/khoang_cach/100m.png}{100~m}
    \caption{Một số mẫu spectrogram theo khoảng cách thu}
    \label{fig:appendix_distance_examples}
\end{figure}

\end{document}
